% =========================================================
% precedentguard_theorems_v0.2_skeleton.tex
% PrecedentGuard — Complete Proof Skeleton for v0.2 Theorems
% AAAI 2027 Submission Working Document
% Created: 2026-06-30 evening (Day 0)
% Proofs due: Day 1 (July 1) EOD
%
% Contents:
%   Theorem 1: Directional Intervention Sensitivity Bound
%   Theorem 2: Population Double-Sided Risk Bound
%   Theorem 3: Finite-Sample Double-Sided Certificate
%   Proposition 1: TV-based Indistinguishability Lower Bound
%   Corollary 1: Authenticity is not Semantic Information
%
% Status markers:
%   [FILL: ...] = you must complete this content
%   [VERIFY: ...] = numerical / definitional choice to confirm
%   [REF: ...] = citation to add at this point
% =========================================================

\documentclass[11pt]{article}
\usepackage[margin=1in]{geometry}
\usepackage{amsmath, amssymb, amsthm, mathtools}
\usepackage{algorithm}
\usepackage[noend]{algpseudocode}
\usepackage{hyperref}
\usepackage{xcolor}
\usepackage{enumitem}

\newcommand{\TODO}[1]{{\color{red}\textbf{[FILL: #1]}}}
\newcommand{\VERIFY}[1]{{\color{orange}\textbf{[VERIFY: #1]}}}
\newcommand{\REF}[1]{{\color{teal}\textsf{[REF: #1]}}}
\newcommand{\NOTE}[1]{{\color{blue}\textit{[Note: #1]}}}

\theoremstyle{plain}
\newtheorem{theorem}{Theorem}
\newtheorem{lemma}{Lemma}
\newtheorem{proposition}{Proposition}
\newtheorem{corollary}{Corollary}
\theoremstyle{definition}
\newtheorem{definition}{Definition}
\newtheorem{assumption}{Assumption}
\newtheorem{remark}{Remark}

\DeclareMathOperator{\TV}{TV}
\DeclareMathOperator{\KL}{KL}
\DeclareMathOperator{\clip}{clip}
\DeclareMathOperator{\sign}{sign}
\DeclareMathOperator{\attested}{attested}
\DeclareMathOperator{\type}{type}

\title{PrecedentGuard --- Complete Proof Skeleton\\
       \large Theorems 1--3 and Proposition 1 (v0.2 Framework)}
\author{AAAI 2027 Working Document}
\date{2026-06-30}

\begin{document}
\maketitle

\section{Setup and Notation (Aligned with v0.2 \S 3)}

\subsection{Score decomposition}

The PrecedentGuard score on input $x = (I, A, P_x, R_x)$ is
\[
S_{\mathrm{PG}}(x) = B(I, A) + \clip\!\Big(\sum_{e \in P_x} \beta_e \tilde\delta_e^{\mathrm{PG}}
        + \sum_{i \in R_x} w_i \tilde\delta_i^{\mathrm{PG}},\ -\varepsilon_-,\ \varepsilon_+ \Big),
\]
where $B(I, A) \in \mathbb{R}$ is the frozen guard's intent--action base score,
each per-evidence contribution $\tilde\delta_e^{\mathrm{PG}}$ is clipped to a type-specific interval $[-c_{\type(e)}^-, c_{\type(e)}^+]$,
and the directional trust constraint replaces the lower endpoint by $0$ when $\attested(e, x) = 0$:
\[
\tilde\delta_e^{\mathrm{PG}} =
\begin{cases}
\max\!\bigl(0,\ \tilde\delta_e\bigr), & \attested(e,x) = 0,\\[2pt]
\tilde\delta_e, & \attested(e,x) = 1.
\end{cases}
\]
The decision is $\hat Y(x) = \mathbf{1}[S_{\mathrm{PG}}(x) \ge \theta]$.

\subsection{Attack class and budget}

The adversary perturbs at most $m_k$ evidence nodes of type $k$, summarized in a vector budget
$\mathbf m = (m_{\mathrm{mem}}, m_{\mathrm{obs}}, m_{\mathrm{ret}}, m_{\mathrm{tool}}, m_{\mathrm{prec}})$.
We write $\mathcal A(\mathbf m)$ for the set of allowed attacks.
For an attack $a$, let $S(x)$ denote the clean score and $S(a(x))$ the score under attack
with the same $(I, A)$.

\subsection{Class-conditional margins}

For an example $x$ with safety label $Y \in \{0, 1\}$, the clean margin is
\[
M_1(x) = S(x) - \theta \quad (Y = 1, \text{unsafe}),
\qquad
M_0(x) = \theta - S(x) \quad (Y = 0, \text{safe}).
\]

\subsection{Assumptions (A1--A4, v0.2 \S 5.1)}

\begin{assumption}[Execution-graph fidelity]\label{a1}
Every mutable evidence item that enters $S_{\mathrm{PG}}$ is represented by an
instrumented node in the EIG. Hidden channels outside the instrumentation are out of scope.
\end{assumption}

\begin{assumption}[Bounded contribution interface]\label{a2}
Each per-evidence contribution is clipped to $[-c_{\type(e)}^-, c_{\type(e)}^+]$ before aggregation.
The outer aggregate is clipped to $[-\varepsilon_-, \varepsilon_+]$.
\end{assumption}

\begin{assumption}[Attack budget]\label{a3}
The adversary modifies at most $m_k$ evidence nodes of each type $k$, counting insertions, deletions, and
replacements explicitly. A replacement counts as one modified node.
\end{assumption}

\begin{assumption}[Attestation boundary]\label{a4}
For the primary certificate, the adversary cannot forge or improperly obtain a policy attestation.
Authenticated but semantically unauthorized writes remain in the threat model.
Sensitivity to attestation error $\eta$ is treated in \S 5.5 of the main paper.
\end{assumption}

\NOTE{These assumptions are visible deployment conditions, not universal claims.}

% =========================================================
\section{Theorem 1: Directional Intervention Sensitivity Bound}
% =========================================================

\begin{theorem}[Directional Intervention Sensitivity Bound]\label{thm:1}
Under Assumptions~\ref{a1}--\ref{a4}, for every attack $a \in \mathcal A(\mathbf m)$ and every input $x$,
\[
- \rho_-(\mathbf m) \ \le\ S(a(x)) - S(x) \ \le\ \rho_+(\mathbf m),
\]
where, writing $\Delta_k := c_k^- + c_k^+$ and $m_k^{\mathrm{ins,unattested}} \le m_k$ for the
unauthenticated-insertion sub-budget of type $k$,
\[
\rho_+(\mathbf m) = \min\!\Big\{\varepsilon_+,\ \sum_k m_k \Delta_k\Big\},
\]
\[
\rho_-(\mathbf m) = \min\!\Big\{\varepsilon_-,\ \sum_k \big(m_k - m_k^{\mathrm{ins,unattested}}\big)\Delta_k\Big\}.
\]
In particular, an unauthenticated insertion alone cannot decrease the score.
\end{theorem}

\subsection{Proof of Theorem 1}

\begin{proof}
Let $Z(x) = \sum_e \beta_e \tilde\delta_e^{\mathrm{PG}} + \sum_i w_i \tilde\delta_i^{\mathrm{PG}}$ denote
the clean pre-outer-clip aggregate, and $Z(a(x))$ the attacked aggregate.
Let $T \subseteq P_x \cup R_x$ index the modified terms. Write
\[
Z(a(x)) - Z(x) = \sum_{j \in T} \bigl(z_j^{a} - z_j\bigr),
\]
where each $z_j$ is the (clipped, trust-adjusted) contribution of the $j$-th term.

\paragraph{Step 1 (per-term envelope).}
For a modified term of type $k$, the clean and attacked values both lie in $[-c_k^-, c_k^+]$
after type clipping. Hence
\[
-(c_k^- + c_k^+) \ \le\ z_j^{a} - z_j \ \le\ c_k^- + c_k^+ = \Delta_k.
\]

\paragraph{Step 2 (directional refinement for unauthenticated insertions).}
Consider an \emph{unauthenticated} insertion: an attacker inserts a fresh evidence node $e$ such that
$\attested(e, x) = 0$. Before the directional trust rule, this term could carry a contribution as low as $-c_k^-$.
After applying $\tilde\delta_e^{\mathrm{PG}} = \max(0, \tilde\delta_e)$, the minimum possible value is $0$.
Since the clean configuration assigns $0$ to a non-existent node by convention, the change is
\[
z_j^{a} - z_j \in [0,\ c_k^+] \qquad \text{(unauthenticated insertion of type }k\text{)}.
\]
Hence such terms contribute nothing to the downward shift.

\paragraph{Step 3 (deletions and replacements remain bidirectional).}
A deletion removes an existing $z_j \in [-c_k^-, c_k^+]$ and sets it to $0$, giving change in $[-c_k^+,\ c_k^-]$;
a replacement keeps the attacker in the full bidirectional envelope $[-\Delta_k,\ \Delta_k]$.
Replacements of \emph{already-attested} evidence retain the bidirectional envelope because the trust rule
permits negative attested contributions.

\paragraph{Step 4 (summation over modified terms).}
Summing per-term envelopes over $T$ yields raw bounds
\[
-r_-(\mathbf m) \ \le\ Z(a(x)) - Z(x) \ \le\ r_+(\mathbf m),
\]
with
\[
r_+(\mathbf m) = \sum_k m_k \Delta_k,
\qquad
r_-(\mathbf m) = \sum_k \big(m_k - m_k^{\mathrm{ins,unattested}}\big)\Delta_k.
\]

\paragraph{Step 5 (effect of the outer clip).}
The outer clip $C(z) = \clip(z,\ -\varepsilon_-,\ \varepsilon_+)$ is $1$-Lipschitz and monotone:
for any $z_1, z_2$, $|C(z_1) - C(z_2)| \le |z_1 - z_2|$, and $C$ preserves order.
Therefore
\[
S(a(x)) - S(x) = C(Z(a(x))) - C(Z(x)) \in \big[-r_-,\ r_+\big].
\]
Additionally, the image of $C$ has diameter $\varepsilon_- + \varepsilon_+$, so the change is also bounded
componentwise by $\min(r_+, \varepsilon_+)$ above and $\min(r_-, \varepsilon_-)$ below.
This yields exactly $\rho_+$ and $\rho_-$ as defined.
\end{proof}

\subsection{Remarks on Theorem 1}

\begin{remark}[Tightness]\label{rem:thm1-tight}
The upward bound is attained by an attacker that performs a single insertion of type $k^\star$ with the
largest $c_{k^\star}^+$, with $\beta_e = 1$. The downward bound, restricted to non-insertion attacks,
is attained by a single replacement of an attested positive term by its negative endpoint.
\end{remark}

\begin{remark}[Outer clip is not a redundant safety net]
If the per-type caps $c_k^{\pm}$ are large, the outer clip $\varepsilon_\pm$ becomes the binding constraint
and gives a uniform bound independent of $\mathbf m$. This is the configuration we recommend whenever
the attacker can choose $\mathbf m$ adaptively.
\end{remark}

\begin{remark}[Limit of the bound]
Theorem~\ref{thm:1} does \emph{not} imply that the underlying frozen guard is Lipschitz on arbitrary text.
Robustness comes from preventing mutable evidence from entering an unbounded monolithic context,
not from a textual smoothness property. \TODO{Add a one-line statement here citing why \S 4.2 of v0.2 explicitly avoids this claim.}
\end{remark}

% =========================================================
\section{Theorem 2: Population Double-Sided Risk Bound}
% =========================================================

\begin{theorem}[Population Double-Sided Risk Bound]\label{thm:2}
Under Theorem~\ref{thm:1},
\[
\sup_{a \in \mathcal A(\mathbf m)} \mathrm{FNR}(a)
   \ \le\ \Pr\!\big(M_1 \le \rho_-(\mathbf m) \mid Y = 1\big),
\]
\[
\sup_{a \in \mathcal A(\mathbf m)} \mathrm{FPR}(a)
   \ \le\ \Pr\!\big(M_0 \le \rho_+(\mathbf m) \mid Y = 0\big).
\]
\end{theorem}

\subsection{Proof of Theorem 2}

\begin{proof}
We show the FNR bound; the FPR bound is symmetric.

Let $x$ have $Y = 1$ and suppose $M_1(x) > \rho_-(\mathbf m)$. Then
\[
S(x) - \theta > \rho_-(\mathbf m).
\]
By Theorem~\ref{thm:1},
\[
S(a(x)) \ge S(x) - \rho_-(\mathbf m) > \theta,
\]
so $\hat Y(a(x)) = 1$ (the attacked classifier still blocks), and $x$ is not a false negative under attack $a$.

Therefore, for any attack $a \in \mathcal A(\mathbf m)$, the false-negative event satisfies
\[
\big\{\hat Y(a(x)) = 0,\ Y = 1\big\} \ \subseteq\ \big\{M_1(x) \le \rho_-(\mathbf m),\ Y = 1\big\}.
\]
Taking conditional probabilities and then supremum over $a \in \mathcal A(\mathbf m)$ gives the FNR claim.

For FPR: if $Y = 0$ and $M_0(x) > \rho_+$, then $\theta - S(x) > \rho_+$, hence
$S(a(x)) \le S(x) + \rho_+ < \theta$, so $\hat Y(a(x)) = 0$ and $x$ is not a false positive.
The same containment argument yields the FPR bound.
\end{proof}

\subsection{Remarks on Theorem 2}

\begin{remark}[Clean errors are absorbed]
A clean false negative has $S(x) < \theta$, i.e.\ $M_1(x) < 0 \le \rho_-$. Such examples are already in
the bounding event $\{M_1 \le \rho_-\}$, so the theorem implies $\mathrm{FNR}^{\mathrm{clean}} \le U_{\mathrm{FN}}$
as a special case.
\end{remark}

\begin{remark}[Avoiding the trivial ``always block'' solution]
A vacuous classifier that blocks everything has $\mathrm{FNR} = 0$ but $\mathrm{FPR} = 1$.
The double-sided structure prevents passing this off as ``certified safe'': the FPR bound
exposes the safe-class margin distribution, which collapses to a point mass at $-\theta$ under always-block.
\TODO{Insert a one-paragraph illustration of how the calibration of $\theta$, $\varepsilon_\pm$, $c_k^\pm$
on the dev set jointly controls both $U_{\mathrm{FN}}$ and $U_{\mathrm{FP}}$.}
\end{remark}

% =========================================================
\section{Theorem 3: Finite-Sample Double-Sided Certificate}
% =========================================================

Let $C_1$ and $C_0$ be \emph{independent} held-out calibration subsets containing $n_1$ unsafe and $n_0$ safe
examples, respectively. For a configuration $\gamma = (\theta, \mathbf c, \varepsilon_-, \varepsilon_+)$, define
\[
\widehat R_{\mathrm{FN}}(\gamma) = \frac{1}{n_1} \sum_{x \in C_1} \mathbf{1}\!\big[M_1(x) \le \rho_-(\gamma)\big],
\quad
\widehat R_{\mathrm{FP}}(\gamma) = \frac{1}{n_0} \sum_{x \in C_0} \mathbf{1}\!\big[M_0(x) \le \rho_+(\gamma)\big].
\]
Let $\Gamma$ be the (finite) grid of configurations under consideration.

\begin{theorem}[Finite-Sample Double-Sided Certificate]\label{thm:3}
With probability at least $1 - \alpha$ over the draw of $C_0 \cup C_1$, simultaneously for all $\gamma \in \Gamma$,
\[
\sup_{a \in \mathcal A(\mathbf m)} \mathrm{FNR}(a)
   \ \le\ \widehat R_{\mathrm{FN}}(\gamma) + \sqrt{\frac{\log(2|\Gamma|/\alpha)}{2 n_1}},
\]
\[
\sup_{a \in \mathcal A(\mathbf m)} \mathrm{FPR}(a)
   \ \le\ \widehat R_{\mathrm{FP}}(\gamma) + \sqrt{\frac{\log(2|\Gamma|/\alpha)}{2 n_0}}.
\]
\end{theorem}

\begin{remark}[Sidedness of the union bound]\label{rem:thm3-sidedness}
We use \emph{one-sided} Hoeffding, giving the multiplier $2$: one upper-tail deviation per class, over $|\Gamma|$ configurations and $|\{\mathrm{FN}, \mathrm{FP}\}| = 2$ classes. Only the upper-tail deviation $R^{\star}(\gamma) \le \widehat R(\gamma) + t$ is required for a valid upper-bound certificate. A conservative two-sided variant would replace $2|\Gamma|$ by $4|\Gamma|$. The two variants differ by
\[
\frac{t_{4|\Gamma|}}{t_{2|\Gamma|}} \;=\; \sqrt{\,\frac{\log(4|\Gamma|/\alpha)}{\log(2|\Gamma|/\alpha)}\,},
\]
which is \emph{independent of the calibration size} $n$ but depends on $|\Gamma|$ and $\alpha$.
For the reference regime $(|\Gamma|, \alpha) = (20, 0.05)$, the two-sided tail is $5.06\%$ larger than the one-sided.
Representative values on the grid $(|\Gamma|, \alpha) \in \{5, 10, 20, 50, 100, 500\} \times \{0.01, 0.05, 0.10\}$
range from $2.97\%$ (at $|\Gamma|=500, \alpha=0.01$) to $7.26\%$ (at $|\Gamma|=5, \alpha=0.10$);
the full table is reproducible from \texttt{scripts/day1\_theorem\_numerical\_example.py}
(saved output: \texttt{scripts/day1\_theorem\_numerical\_example.log}, which also reproduces three externally supplied reference values as a spot-check).
We adopt the tighter one-sided version and cite this remark explicitly in \S 5.4 of the main draft.
\end{remark}

\subsection{Proof of Theorem 3}

\begin{proof}
Fix any $\gamma \in \Gamma$. For an example $x_i \in C_1$, define the Bernoulli random variable
\[
Z_i^{\mathrm{FN}}(\gamma) = \mathbf{1}\!\big[M_1(x_i) \le \rho_-(\gamma)\big].
\]
Since the $x_i$ are i.i.d.\ from the unsafe-class population,
\[
\mathbb{E}\!\big[Z_i^{\mathrm{FN}}(\gamma)\big] = \Pr\!\big(M_1 \le \rho_-(\gamma) \mid Y = 1\big)
   \ =:\ R_{\mathrm{FN}}^{\star}(\gamma).
\]
By Hoeffding's inequality (one-sided),
\[
\Pr\!\Big(R_{\mathrm{FN}}^{\star}(\gamma) - \widehat R_{\mathrm{FN}}(\gamma) > t\Big)
   \ \le\ \exp(-2 n_1 t^2),
\]
\REF{Hoeffding 1963; van der Vaart \emph{Asymptotic Statistics}, Lemma 2.2.7.}

Apply a union bound over $\Gamma$ and over both classes. Let $N = 2 |\Gamma|$
(one-sided Hoeffding per class, over $|\Gamma|$ configurations $\times\ 2$ classes; see Remark~\ref{rem:thm3-sidedness}).
Set
\[
t = \sqrt{\frac{\log(N/\alpha)}{2 n_{\min}}} \quad \text{with } n_{\min} = \min(n_0, n_1),
\]
so that the simultaneous probability of any deviation exceeding $t$ in any of the $N$ tail events is at most $\alpha$.
Then, with probability at least $1 - \alpha$, simultaneously for all $\gamma$ and both classes,
\[
R_{\mathrm{FN}}^{\star}(\gamma) \le \widehat R_{\mathrm{FN}}(\gamma) + \sqrt{\tfrac{\log(N/\alpha)}{2 n_1}},
\qquad
R_{\mathrm{FP}}^{\star}(\gamma) \le \widehat R_{\mathrm{FP}}(\gamma) + \sqrt{\tfrac{\log(N/\alpha)}{2 n_0}}.
\]

Combining with Theorem~\ref{thm:2}:
\[
\sup_{a \in \mathcal A(\mathbf m)} \mathrm{FNR}(a) \le R_{\mathrm{FN}}^{\star}(\gamma)
   \le \widehat R_{\mathrm{FN}}(\gamma) + \sqrt{\tfrac{\log(N/\alpha)}{2 n_1}},
\]
and similarly for FPR.
\end{proof}

\subsection{Remarks on Theorem 3}

\begin{remark}[Calibration--deployment matching]
The bound is distribution-dependent: it estimates $R^{\star}(\gamma)$ on $C_0 \cup C_1$ and inherits validity
only when the deployment risk distribution matches the calibration distribution for the relevant event.
Cross-domain validity must be evaluated empirically (\S 6.5 of v0.2).
\end{remark}

\begin{remark}[Independence of calibration and test]
The proof assumes $C_0$ and $C_1$ are drawn independently from the unsafe and safe class populations,
and that $\gamma$ is selected on a separate development split. \TODO{Confirm in implementation that the
configuration grid $\Gamma$ is fixed \emph{before} touching $C_0 \cup C_1$.}
\end{remark}

\begin{remark}[Conformal-risk-control as a stronger alternative]
A conformal-risk-control variant (\REF{Angelopoulos \& Bates 2023, NeurIPS}) gives marginal guarantees without
finite-grid union-bound losses. The paper defers this extension; \TODO{decide whether to include a
half-paragraph mention in \S 5.4 Remark.}
\end{remark}

% =========================================================
\section{Proposition 1: TV-Indistinguishability Lower Bound}
% =========================================================

For a fixed deployment context, let $Q_1$ and $Q_0$ denote the (attacked) observable distributions of guard
\emph{inputs} for unsafe and safe decisions respectively. Here ``observable'' means the runtime evidence
the guard $h$ receives, including all evidence-channel signals available to $h$.

\begin{proposition}[Double-Sided Indistinguishability Lower Bound]\label{prop:1}
For any (deterministic or randomized) binary guard $h$,
\[
\mathrm{FNR}(h) + \mathrm{FPR}(h) \ \ge\ 1 - \TV(Q_1,\ Q_0),
\]
and hence
\[
\max\!\big\{\mathrm{FNR}(h),\ \mathrm{FPR}(h)\big\} \ \ge\ \frac{1 - \TV(Q_1,\ Q_0)}{2}.
\]
In particular, if $Q_1 = Q_0$, at least one of $\mathrm{FNR}$ or $\mathrm{FPR}$ is at least $1/2$.
\end{proposition}

\subsection{Proof of Proposition 1}

\begin{proof}
The guard $h$ is a binary hypothesis test that, on observing input distributed as either $Q_1$ (unsafe)
or $Q_0$ (safe), outputs a binary decision. The minimum possible sum of type-I and type-II errors over all
(deterministic or randomized) tests between $Q_1$ and $Q_0$ is
\[
\inf_h \Big(\Pr_{Q_1}[h = 0] + \Pr_{Q_0}[h = 1]\Big) = 1 - \TV(Q_1, Q_0)
\]
by Le~Cam's two-point inequality \REF{Le Cam 1986, \emph{Asymptotic Methods in Statistical Decision Theory},
Theorem 2.2}. Hence every binary guard $h$ satisfies
\[
\mathrm{FNR}(h) + \mathrm{FPR}(h) = \Pr_{Q_1}[h = 0] + \Pr_{Q_0}[h = 1] \ge 1 - \TV(Q_1, Q_0).
\]
The maximum-of-two bound follows from $\max(a, b) \ge (a + b)/2$ for non-negative $a, b$.
\end{proof}

\subsection{Remarks on Proposition 1}

\begin{remark}[Role of Proposition 1]
This proposition is a \emph{supporting clarification}, not the paper's primary novelty. Its purpose is to
formalize a precise reason why provenance authenticity (a signature) and semantic authorization (a policy
attestation tied to a safety label) must be separated: if an authenticated channel does not move
$Q_1$ apart from $Q_0$, it cannot improve the classification floor.
\end{remark}

\begin{remark}[Compared to the previously-planned universal impossibility]
The previous plan (v1.x, archived) stated a universal ``no-trust-anchor $\Rightarrow$ no double-sided control''
impossibility. That claim was too broad because a guard may classify perfectly from an entirely different
uncorrupted channel. Proposition~\ref{prop:1} replaces it with the correct, narrower obstruction:
\emph{lack of label-separating information} in the attacker's observation channel.
\end{remark}

% =========================================================
\section{Corollary 1: Authenticity is not Semantic Information}
% =========================================================

\begin{corollary}[Authenticity is not Semantic Information]\label{cor:1}
Let $Z$ be an authenticated variable appended to the guard's observation (a signature, a write-time
provenance tag, etc.). If
\[
\TV\!\big(Q_1^{O, Z},\ Q_0^{O, Z}\big) \ =\ \TV\!\big(Q_1^O,\ Q_0^O\big),
\]
then introducing $Z$ does not improve the lower bound of Proposition~\ref{prop:1}. In particular, a valid
signature whose distribution is independent of the safety label, or that authenticates adversarially chosen
content, cannot break indistinguishability.
\end{corollary}

\begin{proof}
The bound in Proposition~\ref{prop:1} is monotone in the TV distance of the observation channel.
If extending the channel from $O$ to $(O, Z)$ does not increase TV, the bound is unchanged.
\end{proof}

\NOTE{This formalizes why \S 1.2 of v0.2 distinguishes \emph{authenticity}
from \emph{semantic authorization}: a policy attestation is informative about $Y$ only if its conditional
distribution depends on $Y$.}

% =========================================================
\section{Day 1 Task Map (July 1)}
% =========================================================

\subsection{Hour-by-hour schedule}

\begin{enumerate}[label=\textbf{\arabic*.}, leftmargin=*]

\item \textbf{09:00--09:30 --- Compile check.} Compile this file. Fix any LaTeX errors (likely
encoding or package-related). Verify all theorem environments render.

\item \textbf{09:30--11:00 --- Theorem 1 elaboration.} Read Theorem~\ref{thm:1} and its proof.
Resolve every \TODO{...} marker. The most non-trivial step is the case analysis in Steps 2--3:
verify that your codebase's \emph{insertion} convention matches the proof (i.e., a clean ``no node''
contributes $0$).

\item \textbf{11:00--12:00 --- Theorem 2 elaboration.} The proof is short; the work is checking that
$\rho_-(\mathbf m)$ as computed by your implementation matches the value used in the proof.
Run one numerical example end-to-end.

\item \textbf{14:00--16:00 --- Theorem 3 elaboration.} Resolve the \VERIFY{} marker on the constant
$4 |\Gamma|$. Choose one of:
\begin{itemize}
  \item Conservative two-sided ($4|\Gamma|$): simpler narrative, slightly weaker bound.
  \item One-sided ($2|\Gamma|$): tighter, requires explicit ``one-sided'' caveat.
\end{itemize}
Recommendation: use $2|\Gamma|$ and state it explicitly.

\item \textbf{16:00--17:00 --- Proposition 1 elaboration.} Verify Le~Cam citation; this is standard
but the exact reference matters for the rebuttal.

\item \textbf{17:00--18:00 --- Independent review request.} Send this document plus a short summary email
to a colleague with measure-theory background. Ask: ``Are Steps 2--3 of Theorem 1 watertight under
your reading of insertion, deletion, and replacement?''
\end{enumerate}

\subsection{End-of-day checkpoint}

By 18:00 you should have:
\begin{itemize}[noitemsep]
  \item All \TODO{} markers resolved.
  \item The \VERIFY{} marker for Theorem 3 decided.
  \item One numerical example run end-to-end demonstrating the bound.
  \item One independent review request out.
\end{itemize}

% =========================================================
\section{Self-Review Checklist (Before Gate $\alpha$)}
% =========================================================

\begin{enumerate}[label=$\Box$, leftmargin=*, noitemsep]
\item Notation in this skeleton matches v0.2 main draft \S 3 \emph{exactly} (any divergence is corrected before Gate $\alpha$).
\item Assumptions \ref{a1}--\ref{a4} are quoted verbatim from v0.2 \S 5.1.
\item Theorem 1's directional refinement (Step 2) is consistent with the implementation of
$\tilde\delta_e^{\mathrm{PG}} = \max(0, \tilde\delta_e)$ under $\attested = 0$.
\item Theorem 2's containment $\{\hat Y = 0, Y = 1\} \subseteq \{M_1 \le \rho_-\}$ is established
without invoking measure-theoretic regularity (only standard probability).
\item Theorem 3's calibration set $C_0 \cup C_1$ is provably independent of the configuration grid $\Gamma$.
\item Theorem 3's constant ($2|\Gamma|$ vs $4|\Gamma|$) is consistent throughout the paper.
\item Proposition 1's Le~Cam citation is verified (Le Cam 1986 or Tsybakov 2009, chosen consistently).
\item Corollary 1's monotonicity in TV is correctly invoked (no hidden randomization assumption).
\item One numerical worked example (on a toy intervention) confirms the bound is non-vacuous.
\item Independent mathematical review has provided a written response.
\end{enumerate}

\NOTE{All ten boxes must be checked by 2026-07-04 EOD (Gate $\alpha$). Any unchecked box is a
sprint blocker, not a polish item.}

% =========================================================
\section{Failure Modes and Fallbacks}
% =========================================================

\begin{itemize}[noitemsep]
\item \textbf{If Theorem 1's directional refinement is wrong (e.g., insertions can decrease the score under your
actual implementation):} Revert to the symmetric bound $\rho_{\mathrm{sym}} = \min(\varepsilon_-+\varepsilon_+,\ \sum_k m_k \Delta_k)$
and lose the asymmetric narrative.
\item \textbf{If Theorem 3's calibration assumption fails (e.g., $\Gamma$ depends on $C_0 \cup C_1$):}
Apply a holdout split or conformal-risk-control variant. \REF{Angelopoulos \& Bates 2023.}
\item \textbf{If Proposition 1's TV bound is contested by a reviewer:} It is Le Cam's two-point inequality.
Cite Tsybakov 2009 \emph{Introduction to Nonparametric Estimation}, Theorem 2.2, as a textbook source.
\end{itemize}

\end{document}
